% Ngân hàng tự sinh (AI)
\begin{bt}%[8H4N15-1]
% nguon: AI soạn tương tự câu mẫu a8c7865e
% mau: a8c7865e
Cho 4 đoạn thẳng $a = 8\text{ cm}, b = 12\text{ cm}, c = 14\text{ cm}, d = 21\text{ cm}$. Hỏi $a, b$ có tỉ lệ với $c, d$ không? Vì sao?
	\loigiai{
		Ta có:\\
		\[ \dfrac{a}{b} = \dfrac{8}{12} = \dfrac{2}{3} \]\\
		\[ \dfrac{c}{d} = \dfrac{14}{21} = \dfrac{2}{3} \]\\
		Vì $\dfrac{a}{b} = \dfrac{c}{d}$ nên hai đoạn thẳng $a, b$ tỉ lệ với $c, d$.
	}
\end{bt}

\begin{bt}%[8H4N15-1]
% nguon: AI soạn tương tự câu mẫu a8c7865e
% mau: a8c7865e
Cho 4 đoạn thẳng $a = 9\text{ cm}, b = 15\text{ cm}, c = 12\text{ cm}, d = 20\text{ cm}$. Hỏi $a, b$ có tỉ lệ với $c, d$ không? Vì sao?
	\loigiai{
		Ta có:\\
		\[ \dfrac{a}{b} = \dfrac{9}{15} = \dfrac{3}{5} \]\\
		\[ \dfrac{c}{d} = \dfrac{12}{20} = \dfrac{3}{5} \]\\
		Vì $\dfrac{a}{b} = \dfrac{c}{d}$ nên hai đoạn thẳng $a, b$ tỉ lệ với $c, d$.
	}
\end{bt}

\begin{bt}%[8H4H15-1]
% nguon: AI soạn tương tự câu mẫu a8c7865e
% mau: a8c7865e nang_cao
	Cho ba đoạn thẳng $a = 6\text{ cm}$, $b = 9\text{ cm}$ và đoạn thẳng $d = 12\text{ cm}$. Tìm độ dài đoạn thẳng $c$ để hai đoạn thẳng $a, b$ tỉ lệ với $c, d$.
	\loigiai{
		Để hai đoạn thẳng $a, b$ tỉ lệ với $c, d$ thì ta phải có $\dfrac{a}{b} = \dfrac{c}{d}$.\\
		Ta có $\dfrac{a}{b} = \dfrac{6}{9} = \dfrac{2}{3}$.\\
		Suy ra $\dfrac{c}{d} = \dfrac{2}{3}$, tức là $\dfrac{c}{12} = \dfrac{2}{3}$.\\
		Do đó $c = \dfrac{12 \cdot 2}{3} = 8$ (cm).\\
		Thử lại: $\dfrac{a}{b} = \dfrac{6}{9} = \dfrac{2}{3}$ và $\dfrac{c}{d} = \dfrac{8}{12} = \dfrac{2}{3}$, đúng.\\
		Vậy $c = 8\text{ cm}$.
	}
\end{bt}

\begin{bt}%[8H4N15-2]
% nguon: AI soạn tương tự câu mẫu b58f8484
% mau: b58f8484
	Cho tam giác $ABC$ có $AB=18\text{ cm}, AC=24\text{ cm}$. Trên cạnh $AB$ lấy điểm $M$ sao cho $AM=6\text{ cm}$. Từ $M$ kẻ đường thẳng song song với $BC$ cắt $AC$ tại $N$.
	\begin{enumerate}
		\item Tính độ dài $AN$.
		\item Nếu $BC=30\text{ cm}$, tính $MN$.
	\end{enumerate}
	\loigiai{
		\begin{enumerate}
			\item Vì $MN \parallel BC$, áp dụng định lí Thalès ta có:\\
			$\dfrac{AM}{AB} = \dfrac{AN}{AC} \Rightarrow \dfrac{6}{18} = \dfrac{AN}{24} \Rightarrow AN = \dfrac{24}{3} = 8\text{ cm}$.\\
			\item Áp dụng hệ quả định lí Thalès:\\
			$\dfrac{MN}{BC} = \dfrac{AM}{AB} \Rightarrow \dfrac{MN}{30} = \dfrac{6}{18} = \dfrac{1}{3} \Rightarrow MN = 10\text{ cm}$.
		\end{enumerate}
	}
\end{bt}

\begin{bt}%[8H4N15-2]
% nguon: AI soạn tương tự câu mẫu b58f8484
% mau: b58f8484
	Cho tam giác $ABC$ có $AB=21\text{ cm}, AC=28\text{ cm}$. Trên cạnh $AB$ lấy điểm $M$ sao cho $AM=9\text{ cm}$. Từ $M$ kẻ đường thẳng song song với $BC$ cắt $AC$ tại $N$.
	\begin{enumerate}
		\item Tính độ dài $AN$.
		\item Nếu $BC=35\text{ cm}$, tính $MN$.
	\end{enumerate}
	\loigiai{
		\begin{enumerate}
			\item Vì $MN \parallel BC$, áp dụng định lí Thalès ta có:\\
			$\dfrac{AM}{AB} = \dfrac{AN}{AC} \Rightarrow \dfrac{9}{21} = \dfrac{AN}{28} \Rightarrow AN = \dfrac{28 \cdot 9}{21} = 12\text{ cm}$.\\
			\item Áp dụng hệ quả định lí Thalès:\\
			$\dfrac{MN}{BC} = \dfrac{AM}{AB} \Rightarrow \dfrac{MN}{35} = \dfrac{9}{21} = \dfrac{3}{7} \Rightarrow MN = 15\text{ cm}$.
		\end{enumerate}
	}
\end{bt}

\begin{bt}%[8H4H15-2]
% nguon: AI soạn tương tự câu mẫu b58f8484
% mau: b58f8484 nang_cao
	Cho tam giác $ABC$ có $AB=24\text{ cm}$, $AC=30\text{ cm}$. Trên cạnh $AB$ lấy điểm $M$ sao cho $AM=8\text{ cm}$. Từ $M$ kẻ đường thẳng song song với $BC$ cắt $AC$ tại $N$. Biết rằng $MN=10\text{ cm}$.
	\begin{enumerate}
		\item Tính độ dài $AN$.
		\item Tính độ dài $BC$.
	\end{enumerate}
	\loigiai{
		\begin{enumerate}
			\item Vì $MN \parallel BC$, áp dụng định lí Thalès ta có:\\
			$\dfrac{AM}{AB} = \dfrac{AN}{AC} \Rightarrow \dfrac{8}{24} = \dfrac{AN}{30} \Rightarrow AN = \dfrac{30 \cdot 8}{24} = 10\text{ cm}$.\\
			\item Áp dụng hệ quả định lí Thalès:\\
			$\dfrac{MN}{BC} = \dfrac{AM}{AB} \Rightarrow \dfrac{10}{BC} = \dfrac{8}{24} = \dfrac{1}{3} \Rightarrow BC = 30\text{ cm}$.
		\end{enumerate}
	}
\end{bt}

\begin{bt}%[8H4N15-3]
% nguon: AI soạn tương tự câu mẫu 02e89fca
% mau: 02e89fca
	Cho tam giác $ABC$, trên cạnh $AB$ lấy điểm $M$, trên cạnh $AC$ lấy điểm $N$ sao cho $AM = 6\text{ cm}, MB = 9\text{ cm}, AN = 8\text{ cm}, NC = 12\text{ cm}$. Chứng minh $MN \parallel BC$.
	\loigiai{
		Xét tam giác $ABC$, ta có:\\
		\[ \dfrac{AM}{MB} = \dfrac{6}{9} = \dfrac{2}{3} \]
		\[ \dfrac{AN}{NC} = \dfrac{8}{12} = \dfrac{2}{3} \]
		Do đó, $\dfrac{AM}{MB} = \dfrac{AN}{NC}$.\\
		Theo định lí Thalès đảo, ta suy ra $MN \parallel BC$.
	}
\end{bt}

\begin{bt}%[8H4N15-3]
% nguon: AI soạn tương tự câu mẫu 02e89fca
% mau: 02e89fca
	Cho tam giác $MNP$, trên cạnh $MN$ lấy điểm $E$, trên cạnh $MP$ lấy điểm $F$ sao cho $ME = 5\text{ cm}, EN = 10\text{ cm}, MF = 7\text{ cm}, FP = 14\text{ cm}$. Chứng minh $EF \parallel NP$.
	\loigiai{
		Xét tam giác $MNP$, ta có:\\
		\[ \dfrac{ME}{EN} = \dfrac{5}{10} = \dfrac{1}{2} \]
		\[ \dfrac{MF}{FP} = \dfrac{7}{14} = \dfrac{1}{2} \]
		Do đó, $\dfrac{ME}{EN} = \dfrac{MF}{FP}$.\\
		Theo định lí Thalès đảo, ta suy ra $EF \parallel NP$.
	}
\end{bt}

\begin{bt}%[8H4H15-3]
% nguon: AI soạn tương tự câu mẫu 02e89fca
% mau: 02e89fca nang_cao
	Cho tam giác $ABC$, trên cạnh $AB$ lấy điểm $M$, trên cạnh $AC$ lấy điểm $N$ sao cho $AM = 6\text{ cm}$, $MB = 9\text{ cm}$, $AN = 8\text{ cm}$, $NC = 12\text{ cm}$. Chứng minh $MN \parallel BC$.
	\loigiai{
		Xét tam giác $ABC$, ta có:\\
		$\dfrac{AM}{MB} = \dfrac{6}{9} = \dfrac{2}{3}$\\
		$\dfrac{AN}{NC} = \dfrac{8}{12} = \dfrac{2}{3}$\\
		Do đó, $\dfrac{AM}{MB} = \dfrac{AN}{NC}$.\\
		Theo định lí Thalès đảo, ta suy ra $MN \parallel BC$.
	}
\end{bt}

\begin{bt}%[8H4H16-1]
% nguon: AI soạn tương tự câu mẫu 4b98e5ff
% mau: 4b98e5ff
	Cho tam giác $ABC$ có $BC=16\text{ cm}$. Gọi $M$, $N$ lần lượt là trung điểm của $AB$, $AC$. Tính độ dài $MN$.
	\loigiai{
		$M$, $N$ lần lượt là trung điểm của $AB$, $AC$ nên $MN$ là đường trung bình của tam giác $ABC$.\\
		Do đó $MN=\dfrac12BC=\dfrac12\cdot 16=8\text{ (cm)}$.
	}
\end{bt}

\begin{bt}%[8H4H16-1]
% nguon: AI soạn tương tự câu mẫu 4b98e5ff
% mau: 4b98e5ff
	Cho tam giác $ABC$ có $BC=24\text{ cm}$. Gọi $M$, $N$ lần lượt là trung điểm của $AB$, $AC$. Tính độ dài $MN$.
	\loigiai{
		$M$, $N$ lần lượt là trung điểm của $AB$, $AC$ nên $MN$ là đường trung bình của tam giác $ABC$.\\
		Do đó $MN=\dfrac12BC=\dfrac12\cdot 24=12\text{ (cm)}$.
	}
\end{bt}

\begin{bt}%[8H4V16-1]
% nguon: AI soạn tương tự câu mẫu 4b98e5ff
% mau: 4b98e5ff nang_cao
	Cho tam giác $ABC$ có $BC=18\text{ cm}$. Gọi $M$, $N$ lần lượt là trung điểm của $AB$, $AC$. Biết $MN=2x+1$ (cm). Tính độ dài $MN$ và tìm $x$.
	\loigiai{
		$M$, $N$ lần lượt là trung điểm của $AB$, $AC$ nên $MN$ là đường trung bình của tam giác $ABC$.\\
		Do đó $MN=\dfrac12BC=\dfrac12\cdot 18=9\text{ (cm)}$.\\
		Vì $MN=2x+1$ nên $2x+1=9$, suy ra $2x=8$, do đó $x=4$.\\
		Vậy $MN=9\text{ cm}$ và $x=4$.
	}
\end{bt}

\begin{bt}%[8H4H16-2]
% nguon: AI soạn tương tự câu mẫu 6f39268e
% mau: 6f39268e
	Cho tam giác $DEF$. Gọi $P$ là trung điểm của $DE$. Qua $P$ kẻ đường thẳng song song với $EF$ cắt $DF$ tại $Q$. Chứng minh $Q$ là trung điểm của $DF$ và $PQ=\dfrac12EF$.
	\loigiai{
		Vì $PQ\parallel EF$ nên theo định lí Thalès: $\dfrac{DQ}{QF}=\dfrac{DP}{PE}=1$ (vì $DP=PE$).\\
		Suy ra $DQ=QF$, tức $Q$ là trung điểm của $DF$.\\
		Khi đó $PQ$ là đường trung bình của tam giác $DEF$, nên $PQ=\dfrac12EF$.
	}
\end{bt}

\begin{bt}%[8H4H16-2]
% nguon: AI soạn tương tự câu mẫu 6f39268e
% mau: 6f39268e
	Cho tam giác $GHK$. Gọi $I$ là trung điểm của $GH$. Qua $I$ kẻ đường thẳng song song với $HK$ cắt $GK$ tại $J$. Chứng minh $J$ là trung điểm của $GK$ và $IJ=\dfrac12HK$.
	\loigiai{
		Vì $IJ\parallel HK$ nên theo định lí Thalès: $\dfrac{GJ}{JK}=\dfrac{GI}{IH}=1$ (vì $GI=IH$).\\
		Suy ra $GJ=JK$, tức $J$ là trung điểm của $GK$.\\
		Khi đó $IJ$ là đường trung bình của tam giác $GHK$, nên $IJ=\dfrac12HK$.
	}
\end{bt}

\begin{bt}%[8H4V16-2]
% nguon: AI soạn tương tự câu mẫu 6f39268e
% mau: 6f39268e nang_cao
	Cho tam giác $ABC$, $M$ là trung điểm của $AB$. Qua $M$ kẻ đường thẳng song song với $BC$ cắt $AC$ tại $N$. Biết $BC=14$ cm và $AC=10$ cm. Chứng minh $N$ là trung điểm của $AC$, tính $MN$ và tính $AN$.
	\loigiai{
		Vì $MN\parallel BC$ nên theo định lí Thalès: $\dfrac{AN}{NC}=\dfrac{AM}{MB}=1$ (vì $AM=MB$).\\
		Suy ra $AN=NC$, tức $N$ là trung điểm của $AC$.\\
		Khi đó $MN$ là đường trung bình của tam giác $ABC$, nên $MN=\dfrac12BC=\dfrac12\cdot 14=7$ (cm).\\
		Vì $N$ là trung điểm của $AC$ nên $AN=\dfrac12AC=\dfrac12\cdot 10=5$ (cm).\\
		Thử lại: $MN=7$ cm $=\dfrac12BC$ và $AN=NC=5$ cm, đúng.
	}
\end{bt}
